\scanpage{5}
\mspage{2)}
\deleted{\unsure{2}}

\[
T\unsure{\subset|\leftarrow}S\supset U=S-T
\]
\hspace*{1em}car. résiduelles toutes $=p>0$.

$S$ noethérien, $U$ \deleted{dense} \deleted{\unsure{schémati|schématiquement}} \deleted{schématiquement} dense

\deleted{\unsure{2)}}
\diagram{\begin{tikzcd}[ampersand replacement=\&,column sep=large,row sep=1ex]
 A\arrow[r,mapsto] \& (A_U,T_p^{\deleted{\unsure{p|\natural|\varphi}}}(A),\varphi)\\
 \substack{\text{schémas}\\\text{abélien}\\\text{sur }S} \& \substack{\text{triples }(A_U,M,\varphi)\\ A_U\ \text{schéma abélien sur }U\\ M\ \text{gpes }p\text{-div sur }T\\ \varphi:T_p^{\deleted{p}}(A_U)\simeq M.}
\end{tikzcd}}

\margin{Pour pleine fidélité \deleted{\unsure{si p ex. $S$ irréd. ou car resp $\neq\ell$}} \deleted{\unsure{Au}} lieu de $U$ dense, \inserted{et $S-U$ de car. $p$} il \deleted{\unsure{ne}} suffit \unsure{(?)|que} que \deleted{\unsure{pour}} tte comp. \deleted{\unsure{irréd.}} \deleted{\unsure{de $S$} \illegible{3} $\unsure{T}$ \unsure{non vide} $U$} \deleted{et $\overline U$ est de complément \unsure{de} $S$} \deleted{\unsure{de chaque}} connexe de $S$ rencontre $U$, \deleted{p.ex. $S$} \deleted{pleine fidélité} (\deleted{pour l'effectivité} sinon contre exemple avec $U$ \inserted{connexe et} non vide \deleted{\unsure{+}} obtenu en attachant deux spectres d'anneaux de val. discrète complets, et $U=$ un des deux pts maximaux \ldots) Y-a-t-il des contre'exemples avec $S$ de \unsure{étages|Noether|dim.} \unsure{non équi|non égaux|$2$} ?}

\uline{1)} \uline{Foncteur pleinement fidèle}

fidèle bien connu, car $A\unsure{-|\longrightarrow}T_{\deleted{\unsure{\ell}}p}^{\deleted{p}}(A)$ est fidèle.

``pl.'' fidèle : il faut prouver que si $u_p:T_p^{\deleted{p}}(A)-T_p^{\deleted{p}}(B)$ est algébrique sur $U$, il l'est sur $S$. Si $U$ est sch. dense dans $S$, cela résulte de $[\,]$ 1.2. Dans le cas général, soit $X_0$ l'adhérence schématique de $U$ dans $X$. Alors on dispose d'un hom $\varphi_0:A_{X_0}\longrightarrow B_{X_0}$, \deleted{\unsure{On veut|Comme}} \inserted{compatible avec $u_p$.}

\deleted{\deleted{Soit $\mathcal J$ l'Idéal qui définit $X_0$ sur $X$, il est nilpotent et \unsure{concentré} sur $T$.} Alors les \unsure{théories|thèses|th.} de Serre-Tate impliquent que l'on peut prolonger la restriction de $u_0$ à $T_{\deleted{\mathrm{red}}0}$ (\inserted{$=T$} muni de la structure induite réduite) \deleted{en} un hom. $A_{T_n}\longrightarrow B_{T_n}$, où $T_n$ est le $n$\textsuperscript{ième} voisinage infinitésimal. Se ramenant au cas $S$ affine, \inserted{$=\operatorname{Spec}\Lambda$} on trouve \deleted{que l'on} un hom $A_{\widehat\Lambda}-B_{\widehat\Lambda}$, $\widehat\Lambda=$ complété de $\Lambda$ pour la top. définie par \deleted{\unsure{le}} l'idéal de $T_0$ dans $S_0$. Comme \unsure{:} l'hom $U\times_{\operatorname{Spec}\widehat\Lambda}\longrightarrow\operatorname{Spec}\Lambda$ est fid plat, on gagne .....} car appliquer alors $[\,]$ 1.1.

% archive sheet number: 4
